\scanpage{53}
\mspage{26}
3. \uline{Étude des sections de} $E_U$ \uline{sur} $\Gamma$.

Soit $U$ un \deleted{modèle élémentaire sur} \inserted{schéma connexe lisse de t.f. (\unsure{géométriquem.})} le corps $K$, d'où $E_U\longrightarrow E_K$, et on se propose d'étudier les sections mod $\Pi_{U/K}$ conjugaison --- plus généralement, on \unsure{veut} en $\unsure{=|\overline m}$ temps regarder les sections $E'_K\longrightarrow E_U$, où $E'_K$ est un sous-groupe ouvert de $E_K$. Ce qui signifie que l'on fait une extension de base finie sur $K$). Si \inserted{$K$ de type fini sur le corps $\mathbb{Q}$ et si} $U$ se plonge dans un schéma \deleted{abélien} \inserted{\unsure{ou} un groupe commutatif \deleted{\unsure{algébrique}} rigide}, l'application
\[
U(K)\longrightarrow \inserted{\text{Classes d'isomorphisme}}\ \text{Sections de }B_U\text{ sur }B_K
\]
i.e. classes de $\Pi_{U/K}$ conjugaison des sections de $\unsure{B_U|E_U}$ sur $\unsure{B_K|E_K}$

est injective. On va examiner d'autres façons \inserted{``géométriques''} de trouver des sections.

Supposons d'abord que $U$ est une courbe algébrique, \deleted{\unsure{qui a $\pi_1$ nul}} \inserted{qui ne soit pas de} type $(0,0)$ ou $(0,1)$ i.e. $\pi_1(\bar U)=\Pi_{U/K}\ne\{0\}$ ; Alors \unsure{on a pour} $\forall\ i\in\widehat U\setminus\bar U$ (``pt à l'infini''), le groupe de lacets $L_i$ fournit un scindage (au moins si $i$ est $/K$ i.e. \unsure{fixé}

\margin{On a choisi un rev. universel $\widetilde U$ de $U$ pour définir $\bar K$, et $E_U,E_{\unsure{K}}$, et $E_U\longrightarrow E_K$}
\margin{(1)\diagram{\begin{tikzcd}[ampersand replacement=\&,row sep=small]
1\\
E_K=\Gamma\arrow[u]\\
E_U\arrow[u,no head]\\
\Pi_{U/K}\arrow[u]\\
1\arrow[u]
\end{tikzcd}}}
\margin{52}
